\section{Introducción}
\label{sec:introduccion}

\subsection{Contextualización del Problema}
\label{subsec:contextualizacion}

El Aeropuerto Internacional Jorge Chávez, operado por Lima Airport Partners (LAP), no
distribuye a sus pasajeros de manera uniforme a lo largo del día. En determinadas franjas
horarias y en puntos concretos del terminal (filtros de seguridad, control migratorio,
mostradores de \emph{check-in}, salas de embarque) se forman concentraciones de personas
notablemente superiores a las del resto del recinto en ese mismo momento. La aglomeración,
por sí sola, no es el problema: lo es cuando la cantidad de personas reunidas en un punto
excede, de forma sostenida, la capacidad de atención dispuesta para ese punto, y aparecen
colas que se alargan, tiempos de espera que comprometen la conexión de un vuelo, o
condiciones de incomodidad en un espacio cerrado. LAP cuenta ya con infraestructura de
videovigilancia de circuito cerrado en el terminal; este proyecto explora si esa misma
infraestructura puede aprovecharse, mediante visión computacional, para producir una
estimación automatizada de esa concentración y de los recorridos entre cámaras, sin
reconocimiento facial. Conviene adelantar desde aquí lo que el resto del informe desarrolla: para
seguir a una persona de una cámara a otra el sistema la individualiza con un identificador
anónimo y la reconoce por la apariencia de su cuerpo, y esa decisión de diseño es la que fija su
régimen jurídico (Secciones~\ref{subsec:restricciones} y~\ref{subsec:analisis-teas}).

El equipo no cuenta, al momento de escribir este informe, con evidencia documentada por LAP (incidentes de
aglomeración registrados, tiempos de espera medidos, quejas de pasajeros contabilizadas) que
cuantifique el problema en el propio terminal.\verificar{incidentes por aglomeración
documentados, tiempos de cola observados y quejas recibidas. Solicitado a LAP como punto C-07.} El proyecto parte, por tanto, de una necesidad
operativa plausible y compartida por la literatura de conteo de multitudes en espacios de alto
tránsito, y no de una cifra propia de incidentes del Aeropuerto Jorge Chávez.

\subsection{Importancia e Impacto}
\label{subsec:importancia-impacto}

Conocer cuándo y dónde se concentran las personas es información operativa de valor
directo para un operador aeroportuario: permite anticipar la apertura de mostradores
adicionales, redirigir el flujo de pasajeros hacia un filtro descargado, o dimensionar al personal
de atención con mayor precisión que una revisión manual y periódica, que cubre un punto a la
vez y no deja una serie homogénea sobre la que puedan compararse dos días, dos zonas o dos
franjas horarias. El impacto del proyecto tiene, además, una dimensión que excede lo
operativo: un sistema que observa continuamente a personas en un espacio de uso público
concesionado activa, por diseño, un examen de proporcionalidad bajo la Ley 29733, y el
análisis de impacto de este informe documenta un riesgo distributivo real: si el sistema
subcuenta sistemáticamente a un grupo de personas, las zonas donde ese grupo se concentra
reciben menos recursos de forma sistemática, un daño indirecto pero real que no depende de
identificar a ninguna persona en particular.

\subsection{Motivación desde la Ingeniería}
\label{subsec:motivacion-ingenieria}

Desde la perspectiva de Ingeniería de Sistemas, el proyecto es un caso de estudio útil por
una razón poco común: la restricción dominante no es de cómputo ni de datos, sino normativa.
La mayoría de los trabajos de conteo de multitudes revisados en el marco teórico optimizan
exactitud bajo el supuesto implícito de que capturar más información por persona siempre
mejora el sistema; este proyecto invierte esa prioridad, porque el principio de proporcionalidad
de la Ley 29733 (art. 7) exige que el tratamiento sea adecuado, relevante y no excesivo respecto
de su finalidad. Diseñar bajo esa restricción obliga a hacer explícitas decisiones que en otros
contextos quedarían implícitas: qué se persiste, por cuánto tiempo, y qué mecanismo de
desempate se usa cuando dos lecturas del sistema son ambiguas. La segunda motivación es de
integración bajo restricciones reales de hardware: el prototipo se ejecuta en CPU, sin GPU
dedicada, en computadoras portátiles universitarias, combinando piezas de procedencia distinta
(un detector de personas preentrenado, un seguimiento adaptado, un descriptor de apariencia
preentrenado, una homografía calibrada a mano, una base de datos local y una consola web) en
un único flujo coherente. El costo de esa restricción es concreto: en un procesador de gama
media sin GPU, analizar varias cámaras con una muestra cada 0,2 segundos va entre cinco y
veinticinco veces más lento que el video, por lo que el sistema ofrece modos de rendimiento que
espacian las muestras según el número de cámaras.\verificar{rango de 5 a 25 veces más lento que
el video: estimación del equipo sobre un procesador i5 sin GPU, no una medición formal con
protocolo documentado.}

\subsection{Objetivo General}
\label{subsec:objetivo-general}

Desarrollar y evaluar un sistema de visión computacional capaz de estimar la
concentración y tracking de personas en distintas zonas del Aeropuerto Internacional Jorge
Chávez a partir de imágenes de cámaras fijas, sin recurrir a reconocimiento facial ni a la
extracción de rasgos biométricos identificables.

\subsection{Objetivos Específicos}
\label{subsec:objetivos-especificos}

\begin{enumerate}
  \item Revisar el estado del arte en detección, localización y conteo de personas en escenas
  congestionadas, e identificar qué enfoque técnico resulta más adecuado para un entorno
  aeroportuario y para las restricciones de privacidad del proyecto.
  \item Determinar el marco de cumplimiento normativo aplicable al tratamiento de imágenes
  de un espacio de uso público bajo gestión concesionada, incluyendo el análisis de
  impacto y los requisitos de privacidad desde el diseño.
  \item Diseñar e implementar el componente de detección, seguimiento y asociación de identidad
  entre cámaras, y documentar la granularidad de salida resultante frente al principio de
  proporcionalidad.
  \item Extender el prototipo con analítica de flujo peatonal (cruces de acceso, ocupación por
  zona, dashboard ejecutivo) sobre una cartografía real del terminal, evaluando su
  desempeño con pruebas automatizadas.
  \item Evaluar el desempeño del sistema mediante métricas cuantitativas y cualitativas,
  indicando en cada caso sobre qué conjunto de datos fue medido cada resultado.
\end{enumerate}
